% !PRINTED 311
\weberpage{311}
%<<CONT
and in the sequence
\[
1,\ A_1,\ A_2,\ 1
\]
there thus occur two changes of sign whenever, of the two quantities $A_1$, $A_2$, one or both are negative, and in these cases the equation accordingly has two imaginary roots. But if $A_1$ and $A_2$ are positive, the decision cannot be made from this theorem alone; it is then necessary to take into consideration also the quantity
\[
B = 9\,a_0\,a_3 - a_1\,a_2
\]
and to examine the sign of
\[
3\,D = A_1\,A_2 - B^2
\]
(\S.\,62).

In the case of the biquadratic equations one has to consider the three expressions:
\[
\begin{aligned}
A_1 &= 3\,a_1^2 - 8\,a_0\,a_2,\\
A_2 &= 4\,a_2^2 - 9\,a_1\,a_3,\\
A_3 &= 3\,a_3^2 - 8\,a_2\,a_4.
\end{aligned}
\]
If the signs are, in order, $-+-$, one has four imaginary roots. If all three signs are positive, then four or two or no imaginary roots may be present; in all other cases one has two or four imaginary roots.

\begin{center}
\S.\,102.

The Jacobi Criterion.
\end{center}

In order to estimate the number of roots that lie between two limits $\alpha$, $\beta$, \gesperrt{Jacobi} derived a criterion from the \gesperrt{Cartesian} theorem. If, namely,
\weq{1}{f(x) = 0}
is the proposed equation of degree $n$, of which it is to be decided how many roots it has between the two limits $\alpha$, $\beta$, then \gesperrt{Jacobi}\footnote{Observatiunculae ad theoriam aequationum algebraicarum pertinentes. Crelle's Journal, Bd. 13. Gesammelte Werke, Bd. III.} sets
\weq{2}{y = \frac{x-\alpha}{\beta-x},\qquad x = \frac{\alpha+\beta y}{1+y}}
and now forms the equation
%CONT>>
